%
% Halaman Abstrak
%
% @author  Andreas Febrian
% @version 1.00
%

\chapter*{Abstrak}

\vspace*{0.2cm}
{
	\setlength{\parindent}{0pt}
	
	\begin{tabular}{@{}l l p{10cm}}
		Nama&: & \penulis \\
		Program Studi&: & \program \\
		Judul&: & \judul \\
		Pembimbing&: & \pembimbing \\
	\end{tabular}

	\bigskip
	\bigskip

	Dokumentasi website ini disusun berdasarkan kebutuhan akan pengembangan sebuah website baru, mengingat website lama telah menggunakan teknologi dan basis data yang sudah tidak relevan. Tujuan dari dokumentasi ini adalah untuk memberikan informasi yang komprehensif mengenai proses pengembangan website baru, mulai dari latar belakang, analisis kebutuhan, tahap perancangan dan implementasi, pengujian, hingga cara penggunaannya. Metode yang digunakan dalam pengembangan dan dokumentasi website ini adalah metode Waterfall, yang melalui tahapan berurutan dan sistematis. Hasil dari proses tersebut adalah sebuah website baru berbasis Django yang dirancang sesuai dengan kebutuhan pengguna. Dokumentasi ini diharapkan dapat menjadi pedoman bagi pengguna maupun pengembang di masa mendatang dalam pengelolaan dan pengembangan website lebih lanjut.	\bigskip
	\bigskip
	Kata kunci:\\
	Django, \textit{website}, \textit{scraping}, \textit{framework}
}

\newpage